% $Header: /home/vedranm/bitbucket/beamer/solutions/conference-talks/conference-ornate-20min.en.tex,v 90e850259b8b 2007/01/28 20:48:30 tantau $

\documentclass{beamer} %[compress]
% alternatively add handout or draft option

%\mode<handout>%{\setbeamercolor{background canvas}{bg=black!5}}
%\usepackage{pgfpages}
%\pgfpagesuselayout{4 on 1}[a4paper,border shrink=5mm,landscape]
% For handout printing, 4 slides on 1 page

\mode<presentation>
{
  \usetheme{Pittsburgh} % plain template Pittsburgh
  \useoutertheme{infolines} % bottom (author, title, place)
  \useoutertheme[subsection=false]{miniframes} % top (sections, frames)
  \usecolortheme{beaver} %blue color whale
	\setbeamertemplate{frametitle}[default][right] % frame titles on the right
  \setbeamercovered{transparent}
  % or whatever (possibly just delete it)
}

\setbeamersize{text margin left=1cm, text margin right=1cm}
%\setbeamertemplate{navigation symbols}{}
\usepackage[english]{babel}
\usepackage{booktabs}
\usepackage{comment}
% or whatever

\usepackage[cp1250]{inputenc}
% or whatever
\usepackage{subfigure} %subfigures

\usepackage{times}
%\usepackage{breqn}
\usepackage{amsmath}
\usepackage[T1]{fontenc}
\def\sym#1{\ifmmode^{#1}\else\(^{#1}\)\fi} %shortcut for Stata tables
% Or whatever. Note that the encoding and the font should match. If T1
% does not look nice, try deleting the line with the fontenc.

%\logo{\includegraphics[height=0.5cm]{logo.jpg}} % logo on every page



\title[Lecture 3: Literature Search] % (optional, use only with long paper titles)
{Lecture 3: Literature Search}

\author[T.\ Havranek] % (optional, use only with lots of authors)
{ Tomas Havranek
}
% - Give the names in the same order as the appear in the paper.
% - Use the \inst{?} command only if the authors have different
%   affiliation.

\institute[Prague] % (optional, but mostly needed)
{
  	Charles University, CEPR, METRICS

  %\inst%
  %Economic Research Department\\
  %Czech National Bank
      %\inst{1}%
  %Research Department\\
  %Czech National Bank
  %\and
  %\inst{2}%
  %Institute of Economic Studies\\
  %Charles University, Prague
  
%\pgfdeclareimage[height=1cm]{logo}{logo.jpg} % logo only on title page
%\pgfuseimage{logo}
}% - Use the \inst command only if there are several affiliations.
% - Keep it simple, no one is interested in your street address.

\date[Meta in Econ and Finance] % (optional, should be abbreviation of conference name)
{Research Synthesis in Economics and Finance, \\ University of Canterbury\\[2ex]
  February 24, 2025
}%\\[4mm]% - Either use conference name or its abbreviation.
% - Not really informative to the audience, more for people (including
%   yourself) who are reading the slides online


%\subject{Intertemporal Substitution}

\begin{comment}
\AtBeginSection[]
{
  \begin{frame}<beamer>{Outline}
    \tableofcontents[currentsection]%,currentsubsection]
  \end{frame}
}
\end{comment}

\begin{document}

\begin{frame}
  \titlepage
\end{frame}

\begin{comment}

\begin{frame}{Outline}
  \tableofcontents
  % You might wish to add the option [pausesections]
\end{frame}

\begin{frame}{Inspiration}
\end{frame}

\end{comment}

%\section{Summary}
%\subsection{}

%\section{EIS}
%\subsection{}


\begin{frame}{Choosing your topic}

Prior to starting a meta-analysis, you should have at least a basic understanding of the field. Plus:

\bigskip
  
	\begin{enumerate}
		\item Policy interest
		\item Missing publication bias analysis (in previous meta-analyses)
		\item Missing heterogeneity analysis
		\item Missing focus on economics
		\item New methods
	\end{enumerate}
\end{frame}



\begin{frame}{Inspiration}
\begin{figure}
	\centering
		\includegraphics[width=1.00\textwidth]{meta.png}
\end{figure}
\end{frame}


\begin{frame}{Motivation example 1}
\begin{figure}
	\centering
		\includegraphics[width=0.9\textwidth]{motivation.png}
\end{figure}
\end{frame}


\begin{frame}{Motivation example 2}
\begin{figure}
	\centering
		\includegraphics[width=0.9\textwidth]{motivation7.png}
\end{figure}
\end{frame}


\begin{frame}{Motivation example 3}
\begin{figure}
	\centering
		\includegraphics[width=0.9\textwidth]{motivation3.png}
\end{figure}
\end{frame}


\begin{frame}{Motivation example 4}
\begin{figure}
	\centering
		\includegraphics[width=1.00\textwidth]{motivation5.png}
\end{figure}
\end{frame}


\begin{frame}{Motivation example 5}
\begin{figure}
	\centering
		\includegraphics[width=0.9\textwidth]{motivation4.png}
\end{figure}
\end{frame}


\begin{frame}{Motivation example 6}
\begin{figure}
	\centering
		\includegraphics[width=0.9\textwidth]{motivation2.png}
\end{figure}
\end{frame}

\begin{frame}{Google Scholar}
\textbf{Advantages:}
\begin{itemize}
	\item Full text, not just the abstract, title, keywords
	\item Inclusive
	\item Up-to-date
\end{itemize}

\bigskip

\textbf{Disadvantages:}
\begin{itemize}
	\item Algorithm changes in time (replicability issues)
	\item Too many results
	\item Sometimes relevance unclear
\end{itemize}
\end{frame}

\begin{frame}{Search steps}
\begin{enumerate}
	\item Find 5 most relevant studies (AI can help)
	\item Design a Scholar query that shows the 5 studies near the top
	\item Go through first 500 hits
	\item Read abstracts
	\item Download potentially useful studies
	\item Skim them
	\item Repeat the procedure for the most recent studies (last 3 years)
\end{enumerate}
\end{frame}


\begin{frame}{Search query 1}
\begin{figure}
	\centering
		\includegraphics[width=1\textwidth]{search1.png}
\end{figure}
\end{frame}


\begin{frame}{Search query 2}
\begin{figure}
	\centering
		\includegraphics[width=1\textwidth]{search2.png}
\end{figure}
\end{frame}


\begin{frame}{Primary studies}
\begin{figure}
	\centering
		\includegraphics[width=1\textwidth]{studies.png}
\end{figure}
\end{frame}


\begin{frame}{Snowballing}
\begin{itemize}
	\item No query is perfect
	\item Idea: cover studies not identified by the query but frequently mentioned in the literature
	\item For each primary study, download the list of references
	\item Inspect the 50 studies most frequently cited by the primary studies
	\item If you can use some of them, add them to your list
\end{itemize}
\end{frame}


\begin{frame}{Snowballing example}
\begin{figure}
	\centering
		\includegraphics[width=1\textwidth]{snowball.png}
\end{figure}
\end{frame}

\begin{frame}{Minimum requirements for meta-analysis}
For modern meta-analysis techniques to work, you need at least:

\bigskip

\begin{enumerate}
	\item 30 estimates
	\item 10 studies
	\item effect sizes
	\item standard errors
	\item sample sizes
\end{enumerate}

\end{frame}


\begin{frame}{Quality}
\begin{itemize}
	\item Some studies are published in poor journals or use weak identification. What to do?
	\item \textbf{Option 1}: exclude these studies ex ante. Not recommended.
	\item \textbf{Option 2}: include these studies with controls for journal quality and methodology. You can always exclude them as a robustness check.
	\item Note that some fields have standardized quality benchmarks (risk of bias)
\end{itemize}
\end{frame}


\begin{frame}{Unpublished papers}
\begin{itemize}
	\item Try to collect all studies, published and unpublished
	\item If unfeasible, it is possible to focus just on published studies (fewer typos, peer-review)
	\item Including unpublished papers unlikely to alleviate publication bias
	\item Still too many papers? Recruit a co-author or use a random subset (last resort)
\end{itemize}
\end{frame}


\begin{frame}{PRISMA: document your literature search}
\begin{figure}
	\centering
		\includegraphics[width=0.73\textwidth]{prisma.png}
\end{figure}
\end{frame}



\end{document}


